
% ════════════════════════════════════════════
% 模型建立（5.1.3.模型建立.tex，三级子块）
% ════════════════════════════════════════════
\subsubsection{模型建立}

针对问题一，本节将从能量平衡与储能状态递推的角度建立微网日购电计划的线性规划模型。具体过程如下。

\textbf{步骤1：决策变量与目标函数}

定义模型的决策变量。设第 $k$ 个间隔的购电功率为 $g_k$，储能充电功率为 $c_k$，放电功率为 $d_k$，储电量为 $E_k$，其中 $k=0,1,\dots,143$，$E_k$ 为第 $k$ 个间隔起始时刻的储电量。微网全天购电费用为目标函数，即各间隔购电能量与电价乘积之和：
\begin{equation}
    \min \; Z=\sum_{k=0}^{143} g_k\,p_k\,\Delta t
\end{equation}
其中，$g_k\geq 0$ 为购电功率，$p_k$ 为第 $k$ 个间隔的交易电价，$\Delta t$ 为时间步长。该目标函数线性地刻画了微网的经济诉求。

\textbf{步骤2：功率平衡约束}

微网在每个间隔内须满足小区负载的用电需求。由功率守恒，购电功率、光伏功率与储能放电功率之和须不低于负载功率与储能充电功率之和，即允许在光伏过剩时弃光：
\begin{equation}
    g_k + S_k + d_k - c_k \geq L_k,\quad k=0,1,\dots,143
\end{equation}
其中，$S_k$ 为光伏预测功率，$L_k$ 为小区负载功率，$c_k,d_k$ 分别为充、放电功率。由于光伏发电边际成本为零，最优解会自动优先消纳光伏，只有在光伏与储能放电不足时才从外网购电。

\textbf{步骤3：储能状态递推约束}

储能设备的荷电状态随充放电行为逐间隔更新。考虑到充、放电效率均为 $\eta$，充电时从外部吸收的能量只有 $\eta$ 倍存入电池，放电时向负荷释放单位能量需从电池取出 $1/\eta$ 倍能量，因此递推关系为：
\begin{equation}
    E_{k+1}=E_k+\eta\,c_k\,\Delta t-\frac{1}{\eta}\,d_k\,\Delta t,\quad k=0,1,\dots,143
\end{equation}
其中，$\eta=0.9$ 为储能充放电效率。储能设备须满足容量上下限约束：
\begin{equation}
    E_{\min}\leq E_k\leq E_{\max},\quad E_{\min}=1200,\; E_{\max}=10800
\end{equation}
同时充放电功率受储能变换器额定功率限制：
\begin{equation}
    0\leq c_k\leq P_{\max},\quad 0\leq d_k\leq P_{\max},\quad P_{\max}=5000
\end{equation}

\textbf{步骤4：储电初值与日循环约束}

附录1给定 2025 年 1 月 1 日 0:00 的储电量为 6000 kWh，故储能 0:00 储电量应固定为：
\begin{equation}
    E_0=6000
\end{equation}
同时，问题一要求储能设备在 0:00 与 24:00 的储电量相同，以保证相邻日策略的可重复性，故加入日循环约束：
\begin{equation}
    E_0=E_{144}
\end{equation}
综上所述，问题一的完整线性规划模型为：
\begin{equation}
    \min \; Z=\sum_{k=0}^{143} g_k\,p_k\,\Delta t
\end{equation}
\begin{equation}
    \mathrm{s.t.}\quad g_k+S_k+d_k-c_k\geq L_k
\end{equation}
\begin{equation}
    E_{k+1}=E_k+\eta\,c_k\,\Delta t-\frac{1}{\eta}\,d_k\,\Delta t
\end{equation}
\begin{equation}
    E_{\min}\leq E_k\leq E_{\max},\quad 0\leq c_k,d_k\leq P_{\max},\quad g_k\geq 0
\end{equation}
\begin{equation}
    E_0=6000,\quad E_0=E_{144}
\end{equation}
该模型共有 $4\times 144+1=577$ 个连续变量，其中 $E_0$ 由附录1固定为 6000 kWh；全部约束均为线性等式或不等式，属于标准的线性规划问题\cite{dantzig1963}，可用单纯形法或内点法精确求解，且所得解即全局最优解。

综上所述，本节完成了问题一优化模型的建立，得到了以全天购电费用最小为目标、以功率平衡与储能状态为约束的线性规划模型，接下来将基于此进行求解。
